\centerline{\bf \Huge 6 класс}

\problem{\textbf{1}} \textbf{Критерии:} \\
\textit{7 баллов} $-$ 3 различных правильных разрезания.\\
\textit{4 балла} $-$ 2 различных правильных разрезания.\\
\textit{2 балла} $-$ 1 правильное разрезание.\\
\textit{0-1 балл} $-$ Идея площади (16 клеток => одна часть из 8 клеток)\\
\textit{0 баллов} $-$ Решение неверное, продвижения отсутствуют.\\
\textit{0 баллов} $-$ Решение отсутствует.\\
\textit{0 баллов} $-$ При обоснованных подозрениях на списывание участником из любого источника (Интернет, учебное пособие, ключи олимпиады и т.д.) член жюри имеет право указать на такие фрагменты с помощью восклицательных знаков на полях. \\

\problem{\textbf{2}} \textbf{Критерии суммируются:} \\
\textit{0-4 балла} $-$ Решение А.\\
\textit{0-1 балл} $-$ Составление текстового условия Б.\\
\textit{0-2 балла} $-$ Любое правильное решение Б.\\
\textit{0 баллов} $-$ Любой ответ без объяснений.\\
\textit{0 баллов} $-$ Решение отсутствует.\\
\textit{0 баллов} $-$ При обоснованных подозрениях на списывание участником из любого источника (Интернет, учебное пособие, ключи олимпиады и т.д.) член жюри имеет право указать на такие фрагменты с помощью восклицательных знаков на полях. \\


\problem{\textbf{3}} \textbf{Критерии:} \\
\textit{7 баллов} $-$ Найден верный пример: 3,4,5\\
\textit{0-2 балла} $-$ Верный пример не найден, но имеются продвижения.\\
\textit{0 баллов} $-$ Решение неверное, продвижения отсутствуют.\\
\textit{0 баллов} $-$ Решение отсутствует.\\
\textit{0 баллов} $-$ При обоснованных подозрениях на списывание участником из любого источника (Интернет, учебное пособие, ключи олимпиады и т.д.) член жюри имеет право указать на такие фрагменты с помощью восклицательных знаков на полях. \\

\newpage


\problem{\textbf{4}} \textbf{Критерии:} \\
\textit{7 баллов} $-$ Полное верное решение\\
\textit{5-6 баллов} $-$ Решение в целом верное. Однако оно содержит ошибки, либо пропущенные случаи, не влияющие на логику рассуждений.\\
\textit{0-2 балла} $-$ 2025-9$\cdot$1-90$\cdot$2=1836 цифр на трёхзначные числа\\
\textit{0-2 балла} $-$ 1836:3=612-ое по счёту последнее трёхзначное число\\
\textit{0-2 балла} $-$ Это число 100+612-1=711\\
\textit{0-1 балл} $-$ Нужна последняя цифра числа 711\\
\textit{0 баллов} $-$ Любой ответ без объяснений\\
\textit{0 баллов} $-$ Решение неверное, продвижения отсутствуют.\\
\textit{0 баллов} $-$ Решение отсутствует.\\
\textit{0 баллов} $-$ При обоснованных подозрениях на списывание участником из любого источника (Интернет, учебное пособие, ключи олимпиады и т.д.) член жюри имеет право указать на такие фрагменты с помощью восклицательных знаков на полях.\\

\problem{\textbf{5}} \textbf{Критерии:} \\
\textit{7 баллов} $-$ Полное верное решение\\
\textit{6-7 баллов} $-$ Верное решение. Имеются небольшие недочеты, в целом не влияющие на решение.\\
\textit{5-6 баллов} $-$ Решение в целом верное. Однако оно содержит ошибки, либо пропущенные случаи, не влияющие на логику рассуждений.\\
\textit{4-3 балла} $-$ Верно рассмотрен один из двух (более сложный) существенных случаев. В том случае, когда решение делится на две равноценные части – решение одной из частей. 
\\
\textit{2-3 балла} $-$ Доказаны вспомогательные утверждения, помогающие в решении задачи.
\\
\textit{0-1 балл} $-$ Рассмотрены отдельные важные случаи при отсутствии решения (или при ошибочном решении).\\
\textit{0 баллов} $-$ Любой ответ без объяснений\\
\textit{0 баллов} $-$ Решение неверное, продвижения отсутствуют.\\
\textit{0 баллов} $-$ Решение отсутствует.\\
\textit{0 баллов} $-$ При обоснованных подозрениях на списывание участником из любого источника (Интернет, учебное пособие, ключи олимпиады и т.д.) член жюри имеет право указать на такие фрагменты с помощью восклицательных знаков на полях. \\